\section*{Bab \ref{ch:b1:ecosystem-organization} --- Ekosistem dan Struktur Trofik}

\begin{solution}{exo:b1:ecosystem-organization:1}
\omterm{def:b1:ecosystem-organization:ecosystem}{Ekosistem} kolamnya adalah \omterm{def:b1:ecosystem-organization:ecosystem}{komunitasnya} --- seluruh \omterm{def:b1:populations:population}{populasinya}: alga,
ganggang kolam, siput, larva serangga, katak, ikan, bakteri ---
bersama \omterm{def:b1:ecosystem-organization:ecosystem}{biotopnya}: airnya, suhunya, cahayanya, gas dan mineral
terlarutnya, lumpurnya.
\end{solution}

\begin{solution}{exo:b1:ecosystem-organization:2}
\omterm{def:b1:ecosystem-organization:niche}{Habitat}: tempat sebuah spesies hidup (sebuah hutan cemara). \omterm{def:b1:ecosystem-organization:niche}{Relung}:
caranya hidup di situ --- apa yang dimakannya, kapan, di bagian mana
pohonnya, apa yang ditoleransinya. \omterm{def:b1:ecosystem-organization:niche}{Relung} dasar: jangkau yang dapat
ditempatinya sendirian; \omterm{def:b1:ecosystem-organization:niche}{relung} nyata: apa yang ditinggalkan pesaing
dan pemangsanya baginya.
\end{solution}

\begin{solution}{exo:b1:ecosystem-organization:3}
$20\,800/1.7\times 10^6 = \qty{1.2}{\%}$; herbivoranya mengambil
$3400/8800 = \qty{39}{\%}$ produksi bersihnya.
\end{solution}

\begin{solution}{exo:b1:ecosystem-organization:4}
\omterm{def:b1:ecosystem-organization:trophic}{Produsen}, \omterm{def:b1:ecosystem-organization:trophic}{konsumen} (primer, sekunder, tersier), \omterm{def:b1:ecosystem-organization:trophic}{pengurai}, detritivor;
penguraiannya, yang mengembalikan mineral dari bahan mati kepada
\omterm{def:b1:ecosystem-organization:trophic}{produsennya}.
\end{solution}

\begin{solution}{exo:b1:ecosystem-organization:5}
NPP $= 800\times 17 = \qty{13.6}{MJ/m^2}$; yang dimakan $\qty{2.04}{MJ}$;
yang diasimilasi $\qty{1.02}{MJ}$; kelincinya $\qty{51}{kJ/m^2}$, yakni
\qty{510}{MJ} tiap hektar (\qty{30}{kg} kelinci pada \qty{17}{kJ/g}).
Dayagunanya $51/13\,600 = \qty{0.4}{\%}$.
\end{solution}

\begin{solution}{exo:b1:ecosystem-organization:6}
Biomassa adalah sebuah simpanan: fitoplankton yang dimakan secepat ia
tumbuh tak perlu bertimbun, sehingga tegakan \omterm{def:b1:ecosystem-organization:trophic}{produsennya} dapat lebih
sedikit daripada tegakan hewan berumur lebih panjang yang memakannya.
Energi adalah sebuah alir: tiap arasnya menerima seluruh yang akan
pernah dimilikinya dari aras di bawahnya lalu harus kehilangan
sebagiannya sebagai panas, sehingga alirnya hanya dapat mengecil ke
atas.
\end{solution}

\begin{solution}{exo:b1:ecosystem-organization:7}
$p = 0.5, 0.3, 0.15, 0.05$: $H = -(0.5\ln 0.5 + 0.3\ln 0.3 + 0.15\ln
0.15 + 0.05\ln 0.05) = 0.347 + 0.361 + 0.285 + 0.150 = 1.14$;
kemerataannya $1.14/\ln 4 = 0.82$. Empat pada 25 masing-masing: $H = \ln 4 = 1.39$,
kemerataannya 1.
\end{solution}

\begin{solution}{exo:b1:ecosystem-organization:8}
Bersih $= \qty{3}{mg}$ $\mathrm{O_2}$ tiap liter tiap hari;
respirasinya \qty{1}{mg}; kasar $3 + 1 = \qty{4}{mg}$. Dalam karbon:
kasar \qty{1.5}{mg}, bersih \qty{1.1}{mg} C tiap liter tiap hari.
\end{solution}

\begin{solution}{exo:b1:ecosystem-organization:9}
Tiap pautannya memegang sepersepuluh: sepuluh pautan akan
meninggalkan sepermiliar produksi primernya, terlalu sedikit untuk
memberi makan satu hewan pun pada luas mana pun. Tempat produksinya
rendah (gurun) atau luasnya kecil (pulau), sepersepuluh dari
sepersepuluhnya habis sesudah lebih sedikit pautan; produksi sebuah
hutan yang tinggi menanggung satu atau dua pautan lagi.
\end{solution}

\begin{solution}{exo:b1:ecosystem-organization:10}
Biji-bijian: \qty{1}{t} memerlukan seperenam hektar. Daging sapi:
\qty{1}{t} energi daging sapi memerlukan \qty{10}{t} biji-bijian,
\num{1.7} hektar --- sepuluh kali lahannya. Makan pada aras kedua
memakan biaya energi senilai satu aras; sebuah planet memberi makan
sepuluh kali lebih banyak orang dengan biji-bijian daripada dengan
daging yang diberi makan biji-bijian, dan \omterm{def:b1:ecosystem-organization:trophic}{aras trofik} manusia, yang
dipurata pada seluruh pola makannya, adalah sebuah suku yang besar
dalam \omterm{thm:b1:populations:logistic}{daya dukungnya}.
\end{solution}

\begin{solution}{exo:b1:ecosystem-organization:11}
Tiap meter persegi samudranya adalah sebuah gurun, tetapi ia menutupi
dua pertiga planetnya, sehingga sebuah laju yang rendah dikali sebuah
luas yang sangat besar memberi separuh totalnya. Produksinya dibatasi
haranya, yang tenggelam keluar dari permukaan yang tersinari; di
tempat arusnya membawanya naik kembali ke permukaan (penaikan air
pantai, paparan), produksinya sepuluh kali lebih tinggi, dan di
situ, pada beberapa persen lautnya, ikan dan perikanannya memusat.
\end{solution}

\begin{solution}{exo:b1:ecosystem-organization:12}
Energinya masuk sebagai cahaya, melewati beberapa aras \omterm{def:b1:organism-environment:organism}{organisme}, lalu
pergi seluruhnya sebagai panas: dalam arti itu \omterm{def:b1:ecosystem-organization:ecosystem}{ekosistemnya} memang
mengubah cahaya matahari menjadi panas, dan penguraiannya, yang
merespirasikan segala yang tidak dimakan \omterm{def:b1:ecosystem-organization:trophic}{konsumennya}, adalah tempat
sebagian besarnya pergi. Tetapi materinya tidak habis: karbon,
nitrogen dan fosforus berdaur, yang dikembalikan \omterm{def:b1:ecosystem-organization:trophic}{pengurai} yang sama
kepada \omterm{def:b1:ecosystem-organization:trophic}{produsennya}, sehingga sistemnya bertahan; dan \omterm{def:b1:organism-environment:organism}{organismenya}
bukanlah sebuah hasil sampingan melainkan struktur yang lewatnya
alirnya ditata --- \omterm{def:b1:ecosystem-organization:niche}{relung}, jaring dan piramidanya adalah apa yang
dibangun energinya dalam perjalanannya. Kalimat itu memegang
termodinamikanya lalu melewatkan biologinya.
\end{solution}

\begin{solution}{pb:b1:ecosystem-organization:1}
\textbf{1.} $20\,800 - 12\,000 = \qty{8800}{kcal}$ tiap meter persegi
setahun.
\textbf{2.} \qty{1.2}{\%}. Kehilangannya: cahaya yang tak terserap
(separuh spektrumnya, pantulan, airnya), \omterm{def:b1:flowering-plant-organization:organs}{daun} yang \omterm{def:b1:lipids:fattyacid}{jenuh} pada matahari
penuh atau ternaungi, fotorespirasi, dan respirasi tumbuhannya sendiri
(yang telah dihitung antara yang kasar dan yang bersih).
\textbf{3.} $12\,000/20\,800 = \qty{58}{\%}$.
\textbf{4.} Yang dimakan \qty{3400}{}; langsung ke penguraiannya
$8800 - 3400 = \qty{5400}{kcal}$.
\textbf{5.} $8800/10 = \qty{880}{g/m^2}$: seperti sebuah hutan sedang
atau sebuah padang rumput yang kaya.
\textbf{6.} $1.7\times 10^6 - 20\,800 = \qty{1.68e6}{kcal}$, \qty{98.8}{\%}
darinya: dipantulkan, diteruskan, atau diserap air dan batunya lalu
dipancarkan kembali sebagai panas.
\textbf{7.} Produksinya $1500 - 1100 = \qty{400}{kcal}$; dayaguna
asimilasinya $1500/3400 = \qty{44}{\%}$.
\textbf{8.} $400/1500 = \qty{27}{\%}$.
\textbf{9.} $400/8800 = \qty{4.5}{\%}$ (\qty{16}{\%} bila diperhitungkan
pada apa yang dimakan, \qty{3400}{}).
\textbf{10.} Produksi karnivoranya $380 - 300 = \qty{80}{kcal}$;
dayagunanya $80/400 = \qty{20}{\%}$ (mereka memakan \qty{380}{} dari
\qty{400}{} yang dihasilkan).
\textbf{11.} $6/80 = \qty{7.5}{\%}$.
\textbf{12.} \qty{6}{kcal}; $6/1.7\times 10^6 = \num{3.5e-6}$, beberapa
persejuta.
\textbf{13.} Daging mudah dicerna dan dekat susunannya dengan
pemakannya; bahan tumbuhan berserat dan sebagian besarnya tak tercerna,
sehingga seekor herbivora kehilangan lebih banyak asupannya sebagai
tinja.
\textbf{14.} Produksi bersih yang tak dimakan \qty{5400}{}; tinja
herbivoranya $3400 - 1500 = \qty{1900}{}$; produksi herbivora yang tak
dimakan $400 - 380 = \qty{20}{}$; tinja karnivoranya (asimilasinya tak
diberikan; ambillah yang dimakan dikurangi yang direspirasikan
dikurangi produksinya $= 0$, seluruhnya diasimilasi) dan produksi
karnivora yang tak dimakan $80 - 20 = \qty{60}{}$; produksi karnivora
puncaknya \qty{6}{}: totalnya \qty{7386}{kcal}.
\textbf{15.} $12\,000 + 1100 + 300 + 14 + 7386 = \qty{20800}{kcal}$.
\textbf{16.} Setara dengan produksi kasarnya: mata airnya berada dalam
sebuah keadaan mantap, tanpa menyimpan biomassa dari tahun ke tahun.
\textbf{17.} $7386/20\,800 = \qty{36}{\%}$ (sebagian besar sisanya
adalah respirasi \omterm{def:b1:ecosystem-organization:trophic}{produsennya} sendiri).
\textbf{18.} \omterm{def:b1:ecosystem-organization:trophic}{Produsen} \qty{8800}{}, herbivora \qty{400}{}, karnivora
\qty{80}{}, karnivora puncak \qty{6}{}; dayagunanya \qty{4.5}{\%},
\qty{20}{\%}, \qty{7.5}{\%}.
\textbf{19.} \qty{8000}{}, \qty{400}{}, \qty{100}{}, \qty{15}{kcal/m^2}:
tegak, tiap arasnya sepersepuluh atau kurang dari aras di bawahnya.
\textbf{20.} \omterm{def:b1:ecosystem-organization:trophic}{Produsennya} $8000/8800 = 0.9$ tahun; herbivoranya $400/400 =
1$ tahun; karnivoranya $100/80 = 1.25$ tahun; karnivora puncaknya $15/6 =
2.5$ tahun. Tumbuhannya berganti paling cepat: \omterm{def:b1:body-plans-tissues:tissue}{jaringan} berumur pendek,
yang digembalai sinambung; karnivora puncaknya adalah hewan berumur
panjang yang stoknya diperbarui perlahan.
\textbf{21.} Seekor bangau seberat \qty{2}{kg} sekitar \qty{500}{g}
kering, \qty{5000}{kcal}, lalu memerlukan sekitar itu dalam produksi
tiap tahun untuk menggantikan dirinya dan anaknya:
$5000/6 \approx \qty{800}{m^2}$ mata air --- pada kenyataannya jauh
lebih banyak, sebab
ia juga merespirasi: pada \qty{100}{kcal} sehari, \qty{36500}{kcal}
setahun asupan, yakni \qty{6000}{m^2} aras di bawahnya (produksi
karnivoranya \qty{80}{}), yang menjadi rupa sebuah teritori bangau.
\textbf{22.} Produksi bersih tambahannya tidak dimakan lalu pergi ke
penguraiannya, yang respirasinya berlipat dua; pada malam hari, tanpa
fotosintesis, permintaan oksigennya dapat mengosongkan airnya dari
oksigen lalu membunuh ikannya --- eutrofikasi.
\textbf{23.} Produksi herbivoranya turun menjadi \qty{200}{}:
karnivoranya memperoleh separuh makanannya, produksinya turun menuju
\qty{40}{}, produksi karnivora puncaknya menuju 3, dan bangaunya lebih
sedikit; rumput belutnya, yang lebih sedikit digembalai, menebal dan
lebih banyak darinya pergi tanpa dimakan ke penguraiannya.
\textbf{24.} Energi karnivora puncaknya kecil sekali, tetapi
pengaruhnya ada pada jumlah karnivoranya, yang mengendalikan
herbivoranya, yang mengendalikan rumputnya: menyingkirkannya
membiarkan karnivoranya naik, herbivoranya turun dan rumputnya tumbuh
--- sebuah riam yang menembus jaringnya
(\cref{ch:b1:species-interactions}); menyingkirkan beberapa kilokalori
herbivora tidak mengubah apa pun selain beberapa kilokalori.
\textbf{25.} \omterm{def:b1:ecosystem-organization:trophic}{Produsen} ke herbivora \qty{4.5}{\%} (dari NPP;
\qty{16}{\%} dari apa yang dimakan), herbivora ke karnivora
\qty{20}{\%}, karnivora ke karnivora puncak \qty{7.5}{\%}; cahaya
matahari ke produksi kasar \qty{1.2}{\%}; beberapa persejuta cahaya
mataharinya berakhir dalam karnivora puncaknya.
\end{solution}
